\chapter{Ayrılma Aksiyomları}
\label{ch:separation}

%% ============================================================
\section{Konu}
%% ============================================================

Ayrılma aksiyomları, bir topolojik uzaydaki noktaların ve kapalı kümelerin birbirinden
açık kümeler aracılığıyla ne ölçüde ``ayrılabildiğini'' ölçer.

\begin{table}[h]
\centering
\begin{tabular}{lll}
\toprule
\textbf{Aksiyom} & \textbf{İsim} & \textbf{Koşul} \\
\midrule
T0 & Kolmogorov & $\forall x \neq y$: $\exists U \in \tau$ ayıran \\
T1 & Fréchet & $\forall x \neq y$: iki yönlü ayrım \\
T2 & Hausdorff & $\forall x \neq y$: disjoint $U \ni x$, $V \ni y$ \\
T2.5 & Urysohn & disjoint kapanışlı komşuluklar \\
T3 & Regüler & T1 + nokta/kapalı ayırma \\
T3.5 & Tychonoff & T1 + sürekli fonksiyon ayırma \\
T4 & Normal & T1 + disjoint kapalılar ayırma \\
\bottomrule
\end{tabular}
\caption{Ayrılma aksiyomları}
\end{table}

\noindent\textbf{Sıralama:} T4 $\Rightarrow$ T3.5 $\Rightarrow$ T3 $\Rightarrow$ T2.5
$\Rightarrow$ T2 $\Rightarrow$ T1 $\Rightarrow$ T0.

%% ============================================================
\section{Teoremler}
%% ============================================================

\begin{theorem}[Ayrılma Zinciri]
T4 $\Rightarrow$ T3.5 $\Rightarrow$ T3 $\Rightarrow$ T2.5 $\Rightarrow$ T2 $\Rightarrow$ T1
$\Rightarrow$ T0. Tersi genel olarak doğru değildir.
\end{theorem}

\begin{theorem}[Urysohn Fonksiyon Teoremi]
$X$ normal ise ve $C, D$ disjoint kapalı kümeler ise, $f: X \to [0,1]$ sürekli vardır:
$f|_C \equiv 0$, $f|_D \equiv 1$.
\end{theorem}

\begin{theorem}[Tietze Genişleme]
$X$ normal ise her kapalı $A \subseteq X$ üzerinde $f: A \to [a,b]$ süreklisi
tüm $X$'e sürekli genişletilebilir.
\end{theorem}

\begin{theorem}[Tychonoff Karakterizasyonu]
$X$, T3.5'tir $\iff$ $X$ bir küp $[0,1]^I$'ya homeomorf gömülebilir.
\end{theorem}

\begin{theorem}[Sonlu T1 $\iff$ Ayrık]
Sonlu bir uzayda T1 $\iff$ ayrık topoloji.
\end{theorem}
\begin{proof}[İskelet]
T1 $\Rightarrow$ her tekil küme kapalı $\Rightarrow$ her küme kapalı (sonlu birleşim)
$\Rightarrow$ her küme açık.
\end{proof}

%% ============================================================
\section{Algoritmalar}
%% ============================================================

\begin{algorithm}
\caption{Sonlu T0 Karar Prosedürü}
\begin{algorithmic}[1]
\For{each pair $(x, y)$ with $x \neq y$}
    \If{$\nexists U \in \tau$: $(x \in U \wedge y \notin U)$ or $(y \in U \wedge x \notin U)$}
        \Return \textbf{false}
    \EndIf
\EndFor
\Return \textbf{true}
\end{algorithmic}
\end{algorithm}

\noindent\textbf{Karmaşıklık T0:} $O(|X|^2 \cdot |\tau|)$. \quad
\textbf{Karmaşıklık T2:} $O(|X|^2 \cdot |\tau|^2)$.

%% ============================================================
\section{pytop API}
%% ============================================================

\begin{lstlisting}[language=Python, caption={Ayrılma aksiyom importları}]
from pytop import (
    is_t0, is_t1, is_t2, is_t2_5, is_t3, is_t4,
    is_hausdorff, is_regular, is_normal, is_perfectly_normal,
    separation_chain, analyze_separation,
)
\end{lstlisting}

\begin{table}[h]
\centering
\begin{tabular}{lll}
\toprule
\textbf{Fonksiyon} & \textbf{İmza} & \textbf{Açıklama} \\
\midrule
\texttt{is\_t0} \ldots \texttt{is\_t4} & \texttt{(space)} & Aksiyom kontrolü \\
\texttt{separation\_chain} & \texttt{(space)} & Tüm zincir: \texttt{dict[str, Result]} \\
\texttt{analyze\_separation} & \texttt{(space, property='hausdorff')} & Tek aksiyom sorgu \\
\bottomrule
\end{tabular}
\caption{Ayrılma fonksiyonları}
\end{table}

%% ============================================================
\section{Örnekler}
%% ============================================================

\subsection{Örnek 1: Sierpiński — T0 $\checkmark$, T1 $\times$}

\begin{lstlisting}[language=Python]
from pytop import sierpinski_space, is_t0, is_t1, is_t2
s = sierpinski_space()
print("T0:", is_t0(s).status)
print("T1:", is_t1(s).status)
print("T2:", is_t2(s).status)
\end{lstlisting}
\begin{lstlisting}[caption={Çıktı}]
T0: true
T1: false
T2: false
\end{lstlisting}

\subsection{Örnek 2: Kosonlu $\mathbb{N}$ — T1 $\checkmark$, T2 $\times$}

\begin{lstlisting}[language=Python]
from pytop import naturals_cofinite, is_t1, is_t2
nc = naturals_cofinite()
print("T1:", is_t1(nc).status)
print("T2:", is_t2(nc).status)
\end{lstlisting}
\begin{lstlisting}[caption={Çıktı}]
T1: true
T2: false
\end{lstlisting}

\subsection{Örnek 3: Ayrılma Zinciri — Ayrık}

\begin{lstlisting}[language=Python]
from pytop import discrete_topology, separation_chain
d = discrete_topology(1, 2, 3)
for prop, result in separation_chain(d).items():
    print(f"  {prop:20s}: {result.status}")
\end{lstlisting}
\begin{lstlisting}[caption={Çıktı (kısa)}]
  t0                  : true
  t1                  : true
  hausdorff           : true
  t4                  : true
  perfectly_normal    : true
\end{lstlisting}

%% ============================================================
\section{Alıştırmalar}
%% ============================================================

\subsection*{Kodlama}
\begin{enumerate}[label=\textbf{K\arabic*.}]
  \item \texttt{make\_topology(\{1,2,3\},\{1\},\{2\},\{1,2\})} için \texttt{separation\_chain}
        çalıştırın.
  \item \texttt{finite\_chain\_space(3)} üzerinde en yüksek sağlanan aksiyomu bulun.
  \item \texttt{two\_point\_indiscrete\_space()} üzerinde T0, T1, T2 test edin.
\end{enumerate}

\subsection*{Teori}
\begin{enumerate}[label=\textbf{T\arabic*.}]
  \item T2 $\Rightarrow$ T1 $\Rightarrow$ T0 implicasyonlarını kanıtlayın.
  \item Sonlu uzayda T1 $\iff$ ayrık topoloji olduğunu gösterin.
\end{enumerate}
